% !TeX root = ../../Book1.tex
% 纲要标题：### 22.5 有限层检测、闭包与连续性
% 纲要位置：计划纲要.md，第 760 行。

\section{有限层检测、闭包与连续性}
\label{b1:sec:2205}

闭子群可以通过全部有限商中的像辨认。到有限离散群的同态，其连续性可以通过核是否包含某个基本开子群判定；到赋有逆极限拓扑的群的同态，则可逐个有限坐标检验连续性。本节将这些判别写成公式，并说明商群的拓扑为何仍须保留闭性要求。

% 纲要内容（仅作写作计划，尚非教材正文）：
%
% * 闭子群由其在各有限 Galois 商中的像确定；写出子群闭包的有限商检测公式
% * 以有限域的 Frobenius 生成的稠密循环子群比较 \(\mathbb Z\) 与 \(\widehat{\mathbb Z}\)，说明具有同一不动域的子群为何必须取闭包
% * 连续到有限离散群的同态具有开核；有限层上的表示与相容的连续表示
% * 无限 Galois 对应中的商群结论要求闭正规子群；说明与第四卷连续上同调的联系，不调用其定理
%
% ---
%

\subsection{闭子群与闭包的有限商检测}
\label{b1:subsec:2205:01}
\begin{proposition}\label{b1:prop:closure-finite-quotients}
设 \(L/K\) 为代数 Galois 扩张，\(G=\operatorname{Gal}(L/K)\)，并在 \(G\) 上赋予 Krull 拓扑。令 \(\mathcal N\) 为 \(L/K\) 的全部有限 Galois 子扩张组成的有向集合；对 \(N\in\mathcal N\)，记 \(\pi_N:G\rightarrow\operatorname{Gal}(N/K)\) 为限制映射，\(U_N=\ker\pi_N=\operatorname{Gal}(L/N)\)。则对任意子群 \(H\leq G\)，有
\begin{equation}\label{b1:eq:closure-finite-quotients}
\overline H=\bigcap_{N\in\mathcal N}\pi_N^{-1}\bigl(\pi_N(H)\bigr)
=\bigcap_{N\in\mathcal N}HU_N.
\end{equation}
其中 \(\overline H\) 表示 \(H\) 在 \(G\) 中的闭包。所以两个闭子群相等，当且仅当它们在每个有限 Galois 商中的像相等。
\end{proposition}
\begin{proof}
点 \(\sigma\) 属于 \(\overline H\)，当且仅当每个基本开邻域 \(\sigma U_N\) 与 \(H\) 相交。相交意味着有 \(h\in H\) 满足 \(\pi_N(h)=\pi_N(\sigma)\)，这等价于 \(\sigma\in\pi_N^{-1}\bigl(\pi_N(H)\bigr)\)，也等价于 \(\sigma\in HU_N\)。逐个 \(N\) 求交得到公式。若两个闭子群的有限像相同，两者分别等于公式中的同一个交；反向由取像直接成立。
\end{proof}
把这个公式与\cref{b1:cor:fixed-field-closure}合用，得到一个完整的辨别准则：任意两个子群具有相同不动域，当且仅当它们的闭包相同，也当且仅当它们在全部有限商中的像相同。对子群本身而言，有限信息只确定它的闭包。

\subsection{Frobenius 的幂、闭包与有限子域}
\label{b1:subsec:2205:02}
设 \(q\) 为素数幂，\(F(x)=x^q\)。在 \(G_{\mathbb F_q}=\operatorname{Gal}(\overline{\mathbb F}_q/\mathbb F_q)\cong\widehat{\vphantom{Z}\mathbb Z}\) 中，\(\langle F\rangle\) 对应整数的自然像 \(\mathbb Z\)。\subsecref{b1:subsec:2203:02}已证明这一像稠密而非全群；现在考虑 \(F\) 的幂所生成的子群。

\ProseParagraphBreak
固定正整数 \(d\)，有
\begin{equation}\label{b1:eq:profinite-multiple-kernel}
d\widehat{\vphantom{Z}\mathbb Z}
=\ker\left(\widehat{\vphantom{Z}\mathbb Z}\longrightarrow\mathbb Z/d\mathbb Z\right).
\end{equation}
一个方向由乘 \(d\) 直接得到。反过来，设相容族 \(a\) 在模 \(d\) 下为零。对每个 \(n\)，选择 \(a_{dn}\) 的一个被 \(d\) 整除的整数代表，除以 \(d\) 后取模 \(n\)，记为 \(b_n\)。改变代表只改变 \(n\) 的倍数，故定义不变；原族的相容性说明 \((b_n)_n\) 相容，并满足 \(db=a\)。这证明另一方向。

\ProseParagraphBreak
\(\langle F^d\rangle\) 在第 \(n\) 个商 \(\mathbb Z/n\mathbb Z\) 中的像为
\[
d(\mathbb Z/n\mathbb Z)=\{[dk]_n:k\in\mathbb Z\}
=\frac{d\mathbb Z+n\mathbb Z}{n\mathbb Z}.
\]
这些有限像也正是 \(d\widehat{\vphantom{Z}\mathbb Z}\) 的有限像；后者由\eqref{b1:eq:profinite-multiple-kernel}知是闭子群。因此\eqref{b1:eq:closure-finite-quotients}给出
\(\overline{\langle F^d\rangle}=d\widehat{\vphantom{Z}\mathbb Z}\)。它恰为固定 \(\mathbb F_{q^d}\) 的自同构组成的子群，所以 \(\langle F^d\rangle\) 与其闭包的不动域同为 \(\mathbb F_{q^d}\)。

下表比较这几个子群的拓扑性质与不动域，其中 \(d\geq2\)，\(\mathbb Z\) 表示整数在 \(\widehat{\vphantom{Z}\mathbb Z}\) 中的自然像。\(d\widehat{\vphantom{Z}\mathbb Z}\) 由\eqref{b1:eq:profinite-multiple-kernel}知既闭又开；\(\mathbb Z\) 稠密而非全群，因而不闭，也不可能开。零子群闭；若它开，紧群 \(\widehat{\vphantom{Z}\mathbb Z}\) 就成为有限离散空间，与它无限相矛盾。特别地，\(\mathbb Z\) 与全群虽不相等，却有同一个不动域。
\begin{center}
\begin{tabular}{@{}lccc@{}}
\multicolumn{4}{c}{Frobenius 子群的闭性、开性与不动域}\\
\addlinespace[0.4em]
\toprule
子群 \(H\) & 闭 & 开 & \(\overline{\mathbb F}_q^{\,H}\) \\
\midrule
\(\widehat{\vphantom{Z}\mathbb Z}\) & 是 & 是 & \(\mathbb F_q\) \\
\(d\widehat{\vphantom{Z}\mathbb Z}\) & 是 & 是 & \(\mathbb F_{q^d}\) \\
\(\mathbb Z\) & 否 & 否 & \(\mathbb F_q\) \\
\(\{0\}\) & 是 & 否 & \(\overline{\mathbb F}_q\) \\
\bottomrule
\end{tabular}
\end{center}

\subsection{开核与有限层表示}
\label{b1:subsec:2205:03}
\begin{proposition}\label{b1:prop:finite-continuous-hom}
设 \(L/K\) 为代数 Galois 扩张，\(G=\operatorname{Gal}(L/K)\)，并在 \(G\) 上赋予 Krull 拓扑。对每个有限 Galois 子扩张 \(N/K\)，记 \(G_N=\operatorname{Gal}(N/K)\)，\(\pi_N:G\rightarrow G_N\) 为限制映射。再设 \(D\) 为有限离散群，\(\rho:G\rightarrow D\) 为群同态。以下条件等价：\(\rho\) 连续；\(\ker\rho\) 开；存在某个这样的 \(N\) 及同态 \(\rho_N:G_N\rightarrow D\)，使 \(\rho=\rho_N\pi_N\)。
\end{proposition}
\begin{proof}
连续性使单点 \(\{1\}\) 的原像开。若核开，取 \(U_N\subseteq\ker\rho\)，规定 \(\rho_N(\sigma|_N)=\rho(\sigma)\)。两个提升之商在 \(U_N\) 中，故像相同；限制满射保证公式在全部 \(G_N\) 上有定义，并保持群乘法。反过来，有限层投影连续，有限离散集合之间的任意映射连续，所以经过有限商的同态连续。
\end{proof}
有限层不一定唯一。若某表示经过 \(G_N\)，也经过任何包含 \(N\) 的更大有限 Galois 层；应比较其限制相容性，而不是把选择层的不同误当作表示不同。

\ProseParagraphBreak
取素数 \(\ell\) 和整数 \(d\geq1\)。相容的有限层表示可合成以 \(\ell\) 进整数环为系数的表示。给 \(M_d(\mathbb Z_\ell)\cong\mathbb Z_\ell^{d^2}\) 赋予积拓扑，并给 \(\operatorname{GL}_d(\mathbb Z_\ell)\) 赋予子空间拓扑；这里的矩阵群不取离散拓扑。作为拓扑群，\(\operatorname{GL}_d(\mathbb Z_\ell)\) 可识别为
\[
\varprojlim_r\operatorname{GL}_d(\mathbb Z/\ell^r\mathbb Z).
\]
相容矩阵逐项给出 \(\mathbb Z_\ell\) 中的矩阵，有限层的逆矩阵也相容，从而给出整体逆矩阵。两侧拓扑均由模 \(\ell^r\) 的坐标条件生成，故这一识别还是同胚。具体地，若
\[
V_r=\ker\!\left(\operatorname{GL}_d(\mathbb Z_\ell)
\longrightarrow\operatorname{GL}_d(\mathbb Z/\ell^r\mathbb Z)\right),
\]
则 \(AV_r\)（\(r\geq1\)）构成矩阵 \(A\in\operatorname{GL}_d(\mathbb Z_\ell)\) 的邻域基。若 \(L/K\) 为代数 Galois 扩张，\(G=\operatorname{Gal}(L/K)\) 取 Krull 拓扑，且 \(\rho_r:G\rightarrow\operatorname{GL}_d(\mathbb Z/\ell^r\mathbb Z)\) 均连续、模约化后相容，就得到唯一同态 \(\rho:G\rightarrow\operatorname{GL}_d(\mathbb Z_\ell)\)。目标的基本开集只限制有限个坐标，其原像为有限多个开集的交，故 \(\rho\) 连续；反向将连续的 \(\rho\) 与各坐标投影复合即可。每个 \(\rho_r\) 都可由\cref{b1:prop:finite-continuous-hom}经过某个有限 Galois 商 \(\operatorname{Gal}(N_r/K)\)。这里 \(r\) 记录矩阵系数的精度，\(N_r\) 记录决定这一级表示的域扩张，两者是不同的层次；随着 \(r\) 增大，所需的 \(N_r\) 也可以增大。逐个精度能经过有限商，并不要求存在一个固定有限层同时决定全部精度，因而整体 \(\rho\) 不必有开核或有限像。

\ProseParagraphBreak
以有限域为例，由\cref{b1:thm:finite-field-absolute-galois}，可以把 \(G_{\mathbb F_q}\) 识别为加法群 \(\widehat{\vphantom{Z}\mathbb Z}\)。对相容族 \(a=(a_n)_n\)，只取素数幂模数的坐标，得到
\[
a_\ell=(a_{\ell^r})_{r\geq1}\in\mathbb Z_\ell.
\]
定义二维表示
\[
\rho(a)=\begin{pmatrix}1&a_\ell\\0&1\end{pmatrix}
\in\operatorname{GL}_2(\mathbb Z_\ell).
\]
矩阵乘法将右上角相加，故这是群同态。其模 \(\ell^r\) 约化只依赖 \(a_{\ell^r}\)，因此连续，并经过有限 Galois 层 \(N_r=\mathbb F_{q^{\ell^r}}\)；全部约化相容，所以 \(\rho\) 连续。普通整数 \(k\) 的像是 \(\begin{psmallmatrix}1&k\\0&1\end{psmallmatrix}\)，不同整数在 \(\mathbb Z_\ell\) 中仍不同，故整体像无限。若所有约化都经过同一个有限商，那么两元素在该商中相等时，其矩阵在每个模 \(\ell^r\) 下都相等，从而整体矩阵相等；这会使 \(\rho\) 本身经过该有限商，像也有限，矛盾。因此这一例中没有统一的有限层。其核也不开：若核开，紧性使其陪集只有有限个，而这些陪集与表示的像一一对应，又会迫使像有限。

\subsection{闭正规子群与 Galois 商群}
\label{b1:subsec:2205:04}
设 \(G\) 为拓扑群，\(H\trianglelefteq G\) 为正规子群，\(q:G\rightarrow G/H\) 为自然映射。在抽象商群 \(G/H\) 上定义\term[shang-tuopu]{商拓扑}[quotient topology]：集合 \(V\subseteq G/H\) 开，当且仅当 \(q^{-1}(V)\) 在 \(G\) 中开。本节应用时，\(G=\operatorname{Gal}(L/K)\)，\(H\) 还是闭子群，并令 \(E=L^H\)。

\begin{proposition}\label{b1:prop:galois-quotient-homeomorphism}
设 \(L/K\) 为代数 Galois 扩张，\(G=\operatorname{Gal}(L/K)\)，并在 \(G\) 上赋予 Krull 拓扑。设 \(H\leq G\) 为闭正规子群，\(E=L^H\)。在 \(G/H\) 上赋予商拓扑，在 \(\operatorname{Gal}(E/K)\) 上赋予 Krull 拓扑。则限制映射诱导的群同构
\(G/H\rightarrow\operatorname{Gal}(E/K)\) 是拓扑群同构。
\end{proposition}
\begin{proof}
先看限制同态 \(r:G\rightarrow\operatorname{Gal}(E/K)\)。目标中一个基本邻域规定某个有限 Galois 子扩张 \(F/K\)、\(F\subseteq E\) 上的作用，其原像在 \(G\) 中也是同样规定 \(F\) 上作用的基本开集，所以 \(r\) 连续。写成 \(r=\bar r q\)，由商拓扑定义可知诱导双射 \(\bar r\) 连续。

若 \(C\subseteq G/H\) 闭，则 \(q^{-1}(C)\) 是紧群 \(G\) 的闭子集，因而紧。由\cref{b1:prop:compact-maps}，其连续像 \(r\bigl(q^{-1}(C)\bigr)=\bar r(C)\) 在 Hausdorff 群 \(\operatorname{Gal}(E/K)\) 中紧，故闭。因此 \(\bar r\) 还是闭映射，逆映射连续，所列同构为同胚。
\end{proof}
若 \(H\) 是非闭正规子群，\(G/H\) 的商拓扑就不可能是 Hausdorff：否则单位元单点 \(\{H\}\) 闭，由商映射 \(q\) 的连续性，\(H=q^{-1}(\{H\})\) 也闭，矛盾。以 \(G=\widehat{\vphantom{Z}\mathbb Z}\)、\(H=\mathbb Z\) 的稠密像为例，商群非平凡，商拓扑却只有空集和全空间两个开集。事实上，若商中有非空开集，其原像 \(V\) 是非空开集且满足 \(V+H=V\)。对任意 \(g\in G\)，开集 \(g-V\) 与稠密的 \(H\) 相交，于是 \(g\in V+H=V\)，故 \(V=G\)。

\ProseParagraphBreak
无限 Galois 对应将中间域与闭子群配对，有限中间扩张对应开子群，正规中间扩张对应闭正规子群；相应中间扩张的 Galois 群是原群对该闭正规子群的拓扑商群。更进一步的连续群上同调和算术表示，还需指定系数对象及相应的连续性条件。
